% Examples/sample.tex
% A sample LaTeX document that exercises PSLaTeXML features.
\documentclass[12pt]{article}

\usepackage{amsmath}
\usepackage{amssymb}
\usepackage{graphicx}
\usepackage{hyperref}

\title{PSLaTeXML Sample Document}
\author{Alice Smith \and Bob Jones}
\date{\today}

\begin{document}

\maketitle
\tableofcontents

\begin{abstract}
This document serves as a comprehensive sample for testing the
PSLaTeXML PowerShell module, which converts \LaTeX{} documents to
XML or HTML output.
\end{abstract}

\section{Introduction}
\label{sec:intro}

\LaTeX{} is a high-quality typesetting system.
It is the standard for academic publications in many fields.

This document demonstrates \textbf{bold text}, \textit{italic text},
\texttt{monospace text}, and \textsc{small-caps text}.
We can also use \underline{underlined} text.

The following is a footnote.\footnote{This is the footnote content.}
Here is a citation~\cite{lamport1994} and a cross-reference to
Section~\ref{sec:math}.

\section{Mathematics}
\label{sec:math}

\subsection{Inline and Display Math}

Inline mathematics: $E = mc^2$ and $\int_0^1 x\,dx = \frac{1}{2}$.

Display mathematics:
$$
\sum_{n=1}^{\infty} \frac{1}{n^2} = \frac{\pi^2}{6}
$$

\subsection{Numbered Equations}

\begin{equation}
\label{eq:euler}
  e^{i\pi} + 1 = 0
\end{equation}

\begin{equation*}
  \nabla \cdot \mathbf{E} = \frac{\rho}{\varepsilon_0}
\end{equation*}

\subsection{Aligned Equations}

\begin{align}
  a + b &= c \\
  x     &= y + z
\end{align}

\section{Lists}

\subsection{Unordered List}

\begin{itemize}
  \item First item
  \item Second item with \textbf{bold} text
  \item Third item with $x^2$ math
  \begin{itemize}
    \item Nested item A
    \item Nested item B
  \end{itemize}
\end{itemize}

\subsection{Ordered List}

\begin{enumerate}
  \item Step one
  \item Step two
  \item Step three
\end{enumerate}

\subsection{Description List}

\begin{description}
  \item[LaTeXML] A converter from \LaTeX{} to XML/HTML.
  \item[PSLaTeXML] A PowerShell port of LaTeXML.
\end{description}

\section{Tables}

\begin{table}[h]
  \centering
  \caption{Sample Data Table}
  \label{tab:data}
  \begin{tabular}{|l|c|r|}
    \hline
    \textbf{Name} & \textbf{Type} & \textbf{Value} \\
    \hline
    Alpha   & integer & 1   \\
    Beta    & float   & 2.5 \\
    Gamma   & string  & abc \\
    \hline
  \end{tabular}
\end{table}

\section{Verbatim and Code}

\begin{verbatim}
function Get-HelloWorld {
    Write-Output "Hello, World!"
}
\end{verbatim}

\section{Block Environments}

\begin{quote}
This is a block quotation. It can span multiple lines and is typically
rendered with additional indentation.
\end{quote}

\begin{abstract}
A second abstract-style block (for demonstration purposes).
\end{abstract}

\subsection{Theorems}

\begin{theorem}
For all $n > 0$, the sum $\sum_{k=1}^{n} k = \frac{n(n+1)}{2}$.
\end{theorem}

\begin{proof}
By induction on $n$.
\end{proof}

\section{Hyperlinks}

Visit \url{https://github.com/brucemiller/latexml} for the original Perl
implementation, and \href{https://github.com/esrklncc793/PSLaTeXML}{PSLaTeXML}
for the PowerShell port.

\section{Accents and Special Characters}

\begin{itemize}
  \item Caf\'{e}, na\"{i}ve, r\^{o}le, pi\~{n}ata
  \item \c{C}a va? --- Oui, \c{c}a va.
  \item En-dash: 1--10; Em-dash: yes---no
  \item Opening quotes: ``hello''; Closing: ``world''
\end{itemize}

\end{document}
